While deep neural networks (DNNs) have demonstrated remarkable capabilities, achieving human-like or even superior performance across a wide range of tasks, their robustness is often compromised by their susceptibility to input perturbations \citep{szegedy}. This vulnerability has catalyzed the formal verification community to develop various methodologies, each presenting a unique trade-off between completeness and computational efficiency~\citep{Marabou,Reluplex,deeppoly}. 
This surge in innovation has also led to competitions, such as VNNComp~\citep{VNNcomp}, to systematically evaluate the performance of neural network verification tools. Notable examples include NNenum~\citep{nnenum}, Marabou~\citep{Marabou,Marabou2}, and PyRAT~\citep{pyrat}, as well as frameworks such as MnBAB~\citep{ferrari2022complete} (which builds upon ERAN~\citep{deeppoly} and PRIMA~\citep{prima}) and $\alpha,\beta$-CROWN~\citep{crown,xu2020fast}, based on branch-and-bound methodology \citep{BaB}.

Most of these verification methods address {\em local} robustness: given a DNN, a specific input image, and a small neighborhood around this image, they verify whether all images in the neighborhood yield the same classification. This neighborhood is provided by a maximal perturbation of the input image, often an $L_\infty$-norm constraint. Under this metric, every subpixel of the input image can vary in a very small range, typically $\frac{2}{255}$ (that is, 2 levels of gray/blue/red/green). Although it is not necessarily the most meaningful perturbation, $L_\infty$ is the usual choice because it is perfectly linear. 
%(subpixel perturbations are independent), which is easier to verify.

%- to line break
Crucially, local robustness verification methods are far too computationally intensive for real-time decision-making pipelines. 
Consider an autonomous vehicle processing a live video feed: images of the feed cannot be certified as robust in a few ms on embedded hardware. This prevents the development of safety filters that could skip non-robust images and only consider certified images.

\smallskip

In this paper, we achieve {\em real-time} certification by leveraging {\em global} robustness bounds that hold across the entire input domain rather than restricting analysis to the neighborhood of a fixed input. Our framework follows a two-step procedure:
 
\textit{Global bound analysis.} The first step is performed offline and computes global bounds on the maximum possible shift between output values of different decision classes due to a perturbation, following an approach similar to VHAGaR~\citep{vhagar}. 
%
Specifically, for any two decision classes $C$ and $D$ and perturbation $\varepsilon$, we compute an upper bound 
$\bar{\beta}^\varepsilon_{C,D}$ on 
$\max_{I,I', |I-I'| \leq \epsilon}(x_C - x'_C + x'_D - x_D)$, 
where $x_{C,D},x'_{C,D}$ are the output values of class $C,D$ for images $I$ and $I'$, respectively.  
%
Unlike local robustness methods, which require repeating computationally expensive calls for each input, the global upper bound $\bar{\beta}^\varepsilon_{C,D}$ is computed offline once per network and remains valid across the entire input space. 

\textit{Real-Time Verification.} The second step occurs online and in real time, after the DNN's inference phase. For a given input image $I$, we identify the class $C$ with the highest output value $x_C$, and check whether for every other class $D \neq C$, 
$x_C - x_D > \bar{\beta}^\varepsilon_{C,D}$. 
% If this is the case, then 
Thereby, we can formally certify that the image $I$ is robust to perturbation $\varepsilon$.
% because an $\varepsilon$-perturbed image $I'$ could at most get $x'_D \leq \bar{\beta}^\varepsilon_{C,D}  - (x_C - x_D)  + x'_C < x'_C$, hence $C$ is also the predicted class for image $I'$. 
%
This check requires only a few dozen CPU instructions, which can be completed under $1$ ms on a single embedded CPU core. When a frame fails this verification, one could either skip it or revert to a safer, degraded mode until a trustworthy, robust frame is received.

Our main contributions address the challenge of computing {\em global bounds} $\bar{\beta}^\varepsilon_{C,D}$ for classes $(C,D)$:

\begin{enumerate}

	\item {\em The AE-DNN architecture for tight bounds.} 
	For vanilla DNNs, global bounds are particularly pessimistic (see Section \ref{sec.pca}), because inputs far away from the training dataset need to be accounted for. This implies that {\em no} techniques can ever produce manageable upper bound $\bar{\beta}^\varepsilon_{C,D}$ for vanilla DNNs.
	Indeed, on inputs not {\em In Distribution (ID)}, DNNs typically behave erratically \citep{OOD}.
	%
	We propose to use the AE-DNN architecture (Fig.~\ref{fig:AE-DNN} from 
	Section~\ref{sec.pca}) {\em in exploitation}, as an alternative
	to the original vanilla DNNs.
	The AE-DNN first uses an auto-encoder (AE) learnt from ID samples, projecting the large input space to a small latent space which aligns well with ID samples. 
	%
	It is then projected back to the full space before being input to the original DNN. Non-ID samples will be projected by the AE-DNN to the ID-compatible latent space, so that they do not create artificially large global bounds. We select the size of the latent space to ensure that the accuracies (which by definition only consider ID inputs) of the AE-DNN and DNN are within $1$ \%. In exploitation, the AE-DNN has similar characteristics than the original DNN,  but global bounds are much better, enabling real-time certification for AE-DNNs.
	%
	We experimented with Linear AEs (LAE): While not as {\color{red}powerful} as non-linear AE, analysing them incur no complexity cost, and they do not require any change of architecture. Experimentally, the latent space is $>30$x smaller than the full space, leading to bounds $\bar{\beta}^\varepsilon_{C,D}$ $3$-$15$x smaller than for vanilla DNNs.
%	 Specifically, we use PCA to faithfully represent common inputs from the dataset. OOD inputs may be represented unfaithfully, which is reasonable as the DNN is unlikely to provide a reasonable answer on OOD inputs anyway. For instance, $20$ linear components represent MNIST inputs faithfully: the PCA-DNN does not lose accuracy on the MNIST dataset.


\item {\em Efficient global bound computation.}
We develop the {\em Diff} MILP encoding for the global robustness problem.
In this model, the variables are the values of the perturbed neurons, as well as the differences between the original and the perturbed neuron values (called the {\em diff variables}, introduced in \citet{diff}). We study and encode how the {\em diff variables} evolve after passing through a ReLU (Prop.~\ref{Prop2}); see Section \ref{s.diff}. By comparison, the {\em classical} MILP model \citep{MILP} employed in \citet{vhagar,lipshitz,ITNE} 
computes the {\em diff variables} as the difference between 2 larger values.
Both encodings are {\em asymptotically} exact (for infinite runtime), but the {\em Diff} model is more accurate in practice (iso finite runtime) due to better linear relaxation.
%, with the same number of variables.
Further, from the {\em Diff MILP model}, which is asymptotically exact, we develop two abstract MILP models, which are more efficient but also asymptotically less accurate than {\em Diff}, namely the {\em 2b+1} and the {\em 1b} models.
	


\iffalse
	\item We develop a novel {\em Diff MILP encoding} for the global robustness problem, 
	where the variables are the values of the perturbed neurons, as well as the difference between the original and the perturbed neuron values (called the {\em diff variables}, introduced in \citep{diff}). We study and encode how the {\em diff variables} evolve after passing through a ReLU (Prop.~\ref{Prop2}), see Section \ref{s.diff}. 
	Compared with the {\em classical MILP model} \citep{MILP} employed in 
	\citep{vhagar,lipshitz,ITNE}, which considers the input and the perturbation instead of the {\em diff variable}, we keep the same number (2) of binary variables per ReLU.
	%the results are more accurate for the same runtime.
	%the linear relaxation is much more accurate, as each {\em diff variable} can be bounded after a ReLU as a function of the value of the {\em diff variables} before the ReLU, whereas the linear relaxations of the classical model is extremely inaccurate, as the variables are independent of each other. This was observed in \citep{lipshitz,ITNE}, and constraints encoding
	%a part of the linear relaxation of our "2v" model were added explicitly. 
	%Experimentally, 
	However, our {\em Diff MILP model} is more efficient, one reason being 
	the accuracy of its linear relaxation.
	%The number of variables a priori doubles compared with local robustness, from each neuron value in the perturbed image to each neuron value in the perturbed image {\em and} in the original image, as the original image is no more fixed. 
	%Recall that the worst case complexity of MILP is exponential in the number of binary variables \citep{DivideAndSlide}. 
	%A straightforward MILP model would be to use the  for each of these variables, as in \citep{vhagar,lipshitz}. The main issue with the classical model is that its linear relaxations is extremely inaccurate, as the variables are independent of each other. Instead, we develop another exact MILP model, 
\fi

	%Namely, the {\em 2b+1} model, accurate on the {\em diff variables} but abstract on the perturbed variables; while the {\em 1b} model has a unique binary variable per ReLU, only considering the {\em diff variables}.

\end{enumerate}


In terms of perturbations, we consider conjunctions of $L_\infty$- and $L_1$-norms, which allow us to accurately describe perturbations. For instance, "each subpixel is perturbed by at most $\frac{50}{255}$ ($L_\infty$) and the sum of the absolute values of perturbations over all subpixels is at most $1$" ($L_1$-perturbation). 
%While $L_1$ perturbations are not linear (because of the absolute values), and thus they are seldom used, they can be used as perturbation in the MILP model without incurring any expensive binary variables (only cheap linear variables are necessary) (Section~\ref{s.L1}). 


Experimentally, 
compared with the optimized versions of the classical MILP encoding
(VHAGaR \citep{vhagar} and ITNE \citep{ITNE}),
the {\em Diff MILP model} computes upper bounds
$\bar{\beta}^{\varepsilon}_{C,D}$ with smaller gaps to the lower bound (see Table \ref{table.classical}). 
%Optimizations are namely reducing the number of binary variables (VHAGaR \citep{vhagar}) or adding linear constraints (ITNE \citep{ITNE}). 
The abstractions {\em 2b+1} and {\em 1b} variants offer different trade-offs, reaching better bounds than the {\em Diff MILP} model when the instance is very complex, or when runtime is limited (Table \ref{table.L1}).  
%Further, using PCA reduces the upper bound $\bar{\beta}^{\varepsilon}_{C,D}$ by $3$ to $20$ times.
Overall, using {\em Diff}, {\em 2b+1}, and {\em 1b}  MILP models 
enables the real-time certification of $\geq 70\%$ of fresh images 
on LAE-DNNs used in exploitation, for an $L_1$-perturbation of $1,4,2$ over the MNIST, Fashion-MNIST, and CIFAR-10 datasets, respectively; see Table \ref{table.cert}. The online process only adds  $\approx 0.5$ ms of latency per image, and $\approx 2000$ images/second can be processed per CPU core.

This is a Concept and Feasibility paper. 
Our techniques are tested on several DNNs and datasets, but the DNN size is limited (though concrete examples are tackled). 
Linear AE and Linear and ReLU activation functions are the ones we experimented with. We strongly believe that the methodology is novel, impactful, and not intrinsically limited to these restrictions, which will eventually be lifted.